\documentclass{beamer}
\usepackage{amsmath, amsfonts, amssymb}
\usepackage{graphicx}
\usetheme{Madrid}
\usecolortheme{wolverine}

\title{GED Math: Algebra Basics, Expressions, and Polynomials}

\date{\\GED Preparation Course}

\begin{document}

%------------------------------------------------
% Title Slide
%------------------------------------------------
\begin{frame}
  \titlepage
\end{frame}

%------------------------------------------------
\section{Lesson 1: Algebra Basics}
%------------------------------------------------

\begin{frame}{What is Algebra?}
\begin{itemize}
    \item Algebra uses symbols (letters) to represent numbers.
    \item Symbols are called \textbf{variables}.
    \item Example: $x + 5 = 9$.
\end{itemize}
\end{frame}

\begin{frame}{Key Concepts}
\begin{itemize}
    \item \textbf{Variable:} A symbol (e.g., $x$, $y$) representing an unknown.
    \item \textbf{Constant:} A fixed number (e.g., 3, -7, 12).
    \item \textbf{Coefficient:} A number multiplying a variable.
    \item \textbf{Equation:} A mathematical statement of equality.
\end{itemize}
\end{frame}

\begin{frame}{Worked Examples}
\textbf{Example 1:} Solve $x + 7 = 12$.
\begin{align*}
  x &= 12 - 7 = 5
\end{align*}
\textbf{Example 2:} Solve $-3x = 12$.
\begin{align*}
  x &= \frac{12}{-3} = -4
\end{align*}
\textbf{Example 3:} Solve $3x + 4 = 19$.
\begin{align*}
  x &= 5
\end{align*}
\textbf{Example 4:} Solve $\frac{x}{4} = 3$.
\begin{align*}
  x &= 12
\end{align*}
\textbf{Example 5:} If $5x = 25$, find $x$.
\begin{align*}
  x &= 5
\end{align*}
\end{frame}

%------------------------------------------------
\section{Lesson 2: Algebraic Expressions}
%------------------------------------------------

\begin{frame}{What are Algebraic Expressions?}
\begin{itemize}
    \item An \textbf{expression} combines variables, constants, and operations.
    \item Example: $3x + 2y - 5$.
    \item No equals sign.
\end{itemize}
\end{frame}

\begin{frame}{Worked Examples}
\textbf{Example 1:} Simplify $4x + 3x - 2 = 7x - 2$.

\textbf{Example 2:} Simplify $2(3x + 4) = 6x + 8$.

\textbf{Example 3:} Simplify $-2(4x - 5) = -8x + 10$.

\textbf{Example 4:} Simplify $5x + 2y - 3x + 4y = 2x + 6y$.

\textbf{Example 5:} Evaluate $2x + 3$ when $x = 4$.
\begin{align*}
  &= 2(4) + 3 = 11
\end{align*}
\end{frame}

%------------------------------------------------
\section{Lesson 3: Polynomials}
%------------------------------------------------

\begin{frame}{What is a Polynomial?}
\begin{itemize}
    \item A \textbf{polynomial} has terms with variables raised to whole number powers.
    \item Example: $3x^2 + 2x + 1$.
    \item The \textbf{degree} is the highest power.
\end{itemize}
\end{frame}

\begin{frame}{Worked Examples}
\textbf{Example 1:} Add $(3x + 2) + (4x + 5) = 7x + 7$.

\textbf{Example 2:} Subtract $(6x^2 + 3x) - (2x^2 + 5x) = 4x^2 - 2x$.

\textbf{Example 3:} Multiply $3x(2x + 4) = 6x^2 + 12x$.

\textbf{Example 4:} Expand $(x + 2)(x + 3) = x^2 + 5x + 6$.

\textbf{Example 5:} Factor $x^2 + 5x + 6 = (x + 2)(x + 3)$.
\end{frame}

%------------------------------------------------
\section{Practice Problems}
%------------------------------------------------

\begin{frame}{Lesson 1 Practice: Algebra Basics}
\begin{enumerate}
    \item Solve $x - 9 = 15$.
    \item Solve $5x = 45$.
    \item Solve $x/3 + 4 = 9$.
    \item Solve $2x - 5 = 11$.
    \item Solve $7x + 3 = 31$.
\end{enumerate}
\end{frame}

\begin{frame}{Lesson 2 Practice: Expressions}
\begin{enumerate}
    \item Simplify $2x + 3x + 4$.
    \item Simplify $3(2x + 5)$.
    \item Simplify $4y - 2(3y - 6)$.
    \item Simplify $5x + 4y - x + 2y$.
    \item Evaluate $3x + 2$ when $x = 6$.
\end{enumerate}
\end{frame}

\begin{frame}{Lesson 3 Practice: Polynomials}
\begin{enumerate}
    \item Add $(2x^2 + 3x + 1) + (x^2 + 4x + 5)$.
    \item Subtract $(5x^2 + 6x) - (3x^2 + 2x)$.
    \item Multiply $2x(3x + 1)$.
    \item Expand $(x + 4)(x + 2)$.
    \item Factor $x^2 + 7x + 10$.
\end{enumerate}
\end{frame}

%------------------------------------------------
\section{Answer Key}
%------------------------------------------------

\begin{frame}{Lesson 1 Answers}
\begin{enumerate}
    \item $x = 24$
    \item $x = 9$
    \item $x = 15$
    \item $x = 8$
    \item $x = 4$
\end{enumerate}
\end{frame}

\begin{frame}{Lesson 2 Answers}
\begin{enumerate}
    \item $5x + 4$
    \item $6x + 15$
    \item $4y - 6y + 12 = -2y + 12$
    \item $4x + 6y$
    \item $20$
\end{enumerate}
\end{frame}

\begin{frame}{Lesson 3 Answers}
\begin{enumerate}
    \item $3x^2 + 7x + 6$
    \item $2x^2 + 4x$
    \item $6x^2 + 2x$
    \item $x^2 + 6x + 8$
    \item $(x + 5)(x + 2)$
\end{enumerate}
\end{frame}

\begin{frame}{Summary}
\begin{itemize}
    \item Algebra helps solve for unknowns using symbols.
    \item Simplify expressions by combining like terms and using distribution.
    \item Polynomials follow consistent rules for addition, subtraction, multiplication, and factoring.
\end{itemize}
\end{frame}

\end{document}
